\chapter{Results}
\label{cha:results}

\section{Preprocessing Validation}

\subsection{Lung Segmentation Quality (TotalSegmentator vs Ground Truth)}
% eda/totalseg_test_dice.csv (via eda/sota.ipynb) -- Dice between
% TotalSegmentator's lung mask and ground truth over the ~163-case test
% split. This is an input-quality check, not a target-abnormality result;
% ties back to the lung-segmentation methodology discussion in Chapter 3
% (classical HU-threshold vs LungMask) as a third data point.

\section{Classical Image Processing Baseline}
% image-processing/outputs/results.csv -- HU-threshold + morphology GGO
% detector (dice/iou/sensitivity/precision), 5 real cases only. A
% pre-deep-learning, pre-sparse-coding baseline; state the small sample
% size explicitly rather than implying it's comparable in scale to the
% other baselines.

\section{Sparse-Coding Dictionary Learning (K-SVD Family)}
% Four dictionary-learning variants are supported by main.py's
% --algorithm flag: ksvd (plain, unsupervised), fddl (Fisher
% Discrimination Dictionary Learning), frozen (IncrementalFrozenDictionary
% -- an LC-KSVD2 base dictionary trained on normal patches only, with each
% abnormality added incrementally as a residual dictionary), and lcksvd
% (the original joint LC-KSVD2, trained over all classes at once). Each
% is covered below to the extent it actually has saved results.

\subsection{K-SVD (Unsupervised)}
% No local saved evaluation output found -- state as trained/available
% via main.py --algorithm ksvd but not yet evaluated locally, rather than
% omitting it silently.

\subsection{FDDL (Fisher Discrimination Dictionary Learning)}
% No local saved evaluation output found -- same caveat as K-SVD above.

\subsection{Frozen Dictionary Learning}

\subsubsection{Sparse Dictionary Interpretability}

\paragraph{Sparsity and Atom Utilisation}
\paragraph{Atom Co-occurrence and Clustering}
\paragraph{Dictionary Coherence}
\paragraph{Visual Inspection of Atoms and Reconstructions}
% (eda.ipynb-derived figures/numbers -- analysed unified_frozen.pkl /
% train_sparse_codes_frozen.npz specifically, i.e. this variant, not
% LC-KSVD2; same material already summarised methodologically in
% Chapter 3's Statistical Analysis)

\subsection{LC-KSVD2 (Joint, Single-Shot)}

\subsubsection{Test-Set Classification Performance}
% Real saved results: outputs/results/lcksvd2_svm_test_*/config.json
% (3 runs) -- patch-level AND scan-level classification report +
% confusion matrix. Note the 3 runs don't agree (one run: ~0.77 patch /
% 0.95 scan accuracy; the other two: ~0.60 patch / 0.52 scan accuracy) --
% report this discrepancy explicitly rather than quietly picking the
% best-looking run.

\subsubsection{Whole-Volume Localisation}
% Qualitative results only -- overlay figures / NIfTI exports from
% inference.py's per-voxel abnormality map (no Dice/IoU computed)

\section{nnU-Net Segmentation Results}
% Dice / IoU / precision / recall / specificity, GGO model only
% (fold 3 + fold 4 ensemble; no nodule-model run completed --
%  state this explicitly rather than omitting it silently)

\section{CT-CLIP Results}

\subsection{Linear-Probe (LiPro) Classification Performance}
% Full-dataset run, 2,465 scans (ct_clip_lipro_classification_metrics.json)

\subsection{Zero-Shot Text-Prompt Pilot}
% Small pilot only (outputs/ctclip_rex_eval) -- flag as pilot-scale,
% not a full comparison run

\subsection{Explainability (Grad-CAM)}
% models/ct-clip/rexground_predictions/full/gradcam_eval_metrics.csv --
% per-case dice/iou/attribution_iou/pointing_game/energy_inside_mask
% against ground-truth masks (~20-case pilot, not the full dataset)

\section{Merlin Results}

\subsection{Zero-Shot Classification Performance}
% merlin_classification_metrics.json

\subsection{Zero-Shot Localisation Performance}
% merlin_localization_summary.csv / per_case.csv -- soft similarity maps,
% not a trained segmentation head

\subsection{Explainability (Grad-CAM)}
% Code exists (explainability/gradcam/gradcam3d.py, run_gradcam.py) but no
% computed metrics file was found -- state this as not yet run rather
% than presenting numbers that don't exist

\section{BiomedParse Results}

\subsection{Presence Detection (Classification) Performance}
% biomed_parse_rexgroundingct_classification_metrics.json

\subsection{Segmentation/Localisation Performance}
% biomed_parse_localization_summary.csv / per_case.csv

\section{MedSAM2 Results}
% results/train_6_a_2/{metrics.json, lesion_bboxes.json} -- single-case
% pilot only (dice 0.0 on that case), from a box-prompted, SAM2
% video-tracking-based lesion segmentation pipeline designed for the full
% 163-case test subset (models/med-sam2/notebooks/
% MedSAM2_RexGroundingCT_Inference.ipynb) but not yet run at that scale.
% Report as a proof-of-concept that the pipeline runs end-to-end, not as
% a comparable evaluation result; note the RECIST-style metric set
% (Hausdorff-95, volume error \%, lesion counts) differs from the
% Dice/IoU used elsewhere.

\section{Cross-Model Comparison}

\subsection{Classification Performance Comparison}
% Shared table/plot: LC-KSVD vs CT-CLIP vs Merlin vs BiomedParse
% (model_comparison_summary.json) -- note LC-KSVD's Atelectasis caveat
% (reuses its Consolidation numbers) if kept in the writeup.
% IMPORTANT: the "LC-KSVD" entry here is sourced from
% models/lc-ksvd/lcksvd_metrics.json, which is a set of separate
% per-abnormality binary LC-KSVD2 dictionaries (e.g. "Lung nodule_lcksvd2.pkl",
% n_components=128) with their own train/val AUROC/F1/AP -- a different
% training setup from the unified 3-class dictionary + shared LinearSVC
% reported under LC-KSVD2 (Joint, Single-Shot) above. Do not conflate the
% two when writing this up; state clearly which one is being compared here.

\subsection{Localisation/Segmentation Performance Comparison}
% nnU-Net vs Merlin vs BiomedParse (LC-KSVD excluded -- no Dice output;
% classical image-processing baseline and MedSAM2 excluded from the main
% table too -- too few cases each to be a fair comparison point, but
% worth a one-line callback to their Results sections above)

\subsection{Explainability Comparison}
% Grad-CAM based: dice / iou / attribution_iou / pointing_game /
% energy_inside_mask. Real computed numbers currently exist only for
% CT-CLIP (~20-case pilot); Merlin's explainability code has not yet been
% run to produce comparable numbers, and BiomedParse has no Grad-CAM
% pipeline (it outputs masks directly rather than a classifier +
% attribution map). Scope this subsection to what's actually comparable,
% and state the gap explicitly rather than presenting an uneven comparison
% as if it were complete.

\section{Chapter Summary}
% One paragraph tying the comparison back to the aims stated in Chapter 1
